\documentclass[aspectratio=169,11pt]{beamer}

%% Shared preamble for Programming and Quantitative Methods in Finance Beamer presentations
%% Adapted from SPM-PEM/spmcourse-beamer.sty (English localisation)

\usepackage{xcolor}
\usepackage{amsmath,amssymb,bm,mathtools}
\usepackage{booktabs}
\usepackage{hyperref}
\usepackage[english]{babel}

\setbeamertemplate{footline}[frame number]
\setbeamertemplate{itemize items}[circle]
\setbeamercolor{structure}{fg=red!30!green!60!black}
\setbeamercolor{item}{fg=red!30!green!60!black}
\setbeamercolor{itemize item}{fg=red!30!green!60!black}

\definecolor{slidebg0}{RGB}{255,255,255}
\definecolor{slidebg1}{RGB}{220,220,220}
\definecolor{slidebg2}{RGB}{220,235,255}
\definecolor{slidebg3}{RGB}{220,255,230}
\definecolor{slidebg4}{RGB}{255,235,210}
\definecolor{slidebg5}{RGB}{235,220,255}
\definecolor{slidebg6}{RGB}{255,220,205}

\newcounter{slidebg}
\setcounter{slidebg}{1}
\newcommand{\alternateslidebg}{
  \ifnum\value{slidebg}=1
    \setbeamercolor{background canvas}{bg=slidebg1}\setcounter{slidebg}{2}
  \else\ifnum\value{slidebg}=2
    \setbeamercolor{background canvas}{bg=slidebg2}\setcounter{slidebg}{3}
  \else\ifnum\value{slidebg}=3
    \setbeamercolor{background canvas}{bg=slidebg3}\setcounter{slidebg}{4}
  \else\ifnum\value{slidebg}=4
    \setbeamercolor{background canvas}{bg=slidebg4}\setcounter{slidebg}{5}
  \else\ifnum\value{slidebg}=5
    \setbeamercolor{background canvas}{bg=slidebg5}\setcounter{slidebg}{6}
  \else
    \setbeamercolor{background canvas}{bg=slidebg6}\setcounter{slidebg}{1}
  \fi\fi\fi\fi\fi
}

\ExplSyntaxOn
\NewExpandableDocumentCommand{\tocsectionbg}{m}
  { slidebg\int_eval:n { \int_mod:nn { #1 } { 6 } + 1 } }
\ExplSyntaxOff

\setbeamertemplate{section in toc}{%
  \setbeamercolor{section in toc}{fg=black,bg=\tocsectionbg{\inserttocsectionnumber}}%
  \begin{beamercolorbox}[wd=\textwidth,sep=0.7ex,leftskip=1em,rightskip=1em]{section in toc}
    \inserttocsectionnumber.\enspace\inserttocsection
  \end{beamercolorbox}%
  \vspace{0.3ex}%
}
\newcommand{\newsection}[1]{\begin{frame}[plain]
\centering\vfill
{\usebeamercolor[fg]{structure}\LARGE\bfseries \insertsectionnumber\ #1}
\vfill\end{frame}}
\NewDocumentCommand{\bgsection}{s m}{%
  \alternateslidebg
  \section{#2}%
  \IfBooleanF{#1}{\newsection{#2}}%
}
\newcommand{\titleandoutline}{
  \alternateslidebg
  \frame{\titlepage}
  \begin{frame}{Outline}
  \tableofcontents[hideallsubsections]
  \end{frame}
}

\title{Programming and Quantitative Methods in Finance}
\subtitle{Python half}
\author{Králik Balázs}
\date{}

\begin{document}


\begin{frame}{How the Python classes work}
  \begin{itemize}
    \item The first part introduces the week's concepts through explanations
          and examples.
    \item The second part is devoted to independent in-class problem solving.
    \item The emphasis is practical: students learn to find and fix syntax,
          implementation, and logical errors while working.
    \item Questions and productive collaboration are encouraged, provided that
          other students can work undisturbed.
  \end{itemize}
\end{frame}

\begin{frame}{Python half: 12 classes}
  \small
  \begin{tabular}{ccll}
    \toprule
    Week & Class & Quiz & Main topic \\
    \midrule
    1 & 1  &           & Basics, variables, and strings \\
      & 2  & $\checkmark$ & Numeric types, conversion, booleans, input, and conditionals \\
    2 & 3  &           & Sequence types, lists, and loops \\
      & 4  & $\checkmark$ & Functions \\
    3 & 5  &           & Review \\
      & 6  & $\checkmark$ & Modules, dates, and the \texttt{math} module \\
    4 & 7  &           & Pandas I \\
      & 8  & $\checkmark$ & Pandas II \\
    5 & 9  &           & Data visualization \\
      & 10 & $\checkmark$ & Monte Carlo simulations \\
    6 & 11 &          & Practice and preparation for the midterm \\
      & 12 &          & Midterm \\
    \bottomrule
  \end{tabular}
\end{frame}

\end{document}
